\section{\label{III_2_A}Logique produit vs logique patrimoniale}

La proposition d'une nouvelle fonctionnalité au sein d'un logiciel de gestion des collections impose de concilier deux logiques distinctes : la logique industrielle de l'éditeur logiciel et la logique de préservation et de gestion des institutions patrimoniales. Ainsi, l'objectif pour l'éiteur n'est pas seulement de proposer des fonctionnalités performantes, mais surtout de proposer des fonctionnalités adaptatives à tous les types de collections potentiellement gérées par les isntitutions clientes. 
Le premier obstacle à l'adaptation d'une fonctionnalité fondée sur l'IA comme c'est le cas de l'indexation automatisée réside dans le paramétrage des modèles de visions par ordinateur. Comme nous l'avons évoqué dans notre étude, un pipeline de segmentation (pyscenedetect) et de détection d'objets (Yolov8) nécessite des ajustements fins en fonction du corpus étudié, notamment en ce qui concerne le seuil de confiance ou le seuil de séquencage. La généralisation de ce pipeline à une clientèle large suppose de renoncer aux modèles lourds que nous avons évoqué tels que les visions transformers, fondés sur le deep learning, qui nécessitent des paramétrages locaux lourds. La tendance s'oriente vers des architectures légères mais hautement interactives afin de permettre à l'utilisateur d'être intégré au fonctionnement de l'algorithme (référence?). Objectif : pouvoir garder le contrôle sur le paramétrage.
Le choix de ce type d'architecture pose la question de l'hébergement des données, pour pouvoir garder le contrôle sur celles-ci. En effet, l'éditeur de logiciel de gestion des collections est reponsable des données clients. Skinsoft propose dans la plupart des cas d'héberger les données de ses clients, afin d'éviter d'installer un serveur de données chez chaque client individuellement et pouvoir facilement gérer les problèmes potentiels. La plupart des outils automatiques disponibles et prêt à l'emploi sont hébergés sur des nuages externes (\textit{clouds}). Dans le cas de l'adoption d'un tel outil en nuage, cela suppose que les données clients sont "envoyées" sur un serveur externe appartenant à une grande entreprise, sur lequel l'éditeur de logiciel n'aurait qu'un contrôle limité, créant ainsi une "dépendance". \footcite{zotero-item-584} En tant qu'éditeur, le rôle de Skinsoft, par exemple, serait de s'approprier un algorithme déjà existant et de le développer pour pouvoir l'adapter à ses clients. De fait, les données seraient entièrement hébérgées par l'éditeur permettant ainsi de gérer les potentiels problèmes sur les données entrantes au sein des algorithmes et les données sortantes. 

De plus, au delà des contraintes posées par de nouvelles fonctionnalités automatisées, se pose la question de la pertinence du déploiement de tels outils, en lien avec les besoins profesionnels actuels. En effet, l'éditeur de logiciel est animé par une logique commerciale qui devance les pratiques ancestrales de gestion des collections. La posture de l'éditeur, en prise avec l'innovation pour devancer toute concurrence contraste donc avec les logiques institutionnelles et patrimoniales de ses clients. Dans son guide pour une IA responsable dans le secteur culturel, le Ministère de la Culture rappelle que dans le cadre d'une mise en place de fonctionnalités fondées sur l'IA, il s'agit "d’examiner le besoin à couvrir et de vérifier la capacité des systèmes d’IA à satisfaire ce besoin dans des conditions optimales" ; le cas échéant, de favoriser des alternatives moins consommatrices que l’IA si elles sont suffisantes à répondre au besoin. » \footcite{2025a}
Dans le cas de Skinsoft, si certains clients sont demandeurs d'une automatisation de certaines tâches réalisables au sien d'un système de gestion des collections, d'autres les refusent catégoriquement de peur de perdre le contrôle sur les données existantes et celles qui seront produites. Le pipeline proposé dans notre étude, fondé sur la computer vision est une alternative intéressante puisqu'elle permet de configurer pleinement les tâches automatisées et de garder le contrôle sur ce qui est produit. Les risques associés habituellement aux modèles d'IA tels que les risque d'hallucination, d'erreurs, de biais, seraient amoindris dans le cadre de ce pipeline, puisque les outils seraient intégrés à un logiciel professionnels comme facilitateurs et propositions, mais jamais comme décisionnaires. 
exemple du NLP : nous avons proposé le NLP comme un module que l'utilisateur peut désactiver s'il le souhaite. 
Les professionnels pourraient mettre à profit l'IA pour un double objectif : la préservation et l'accessibilité des collections. \footcite{zotero-item-640}. C'est exactement le cas pour l'indexation automatique.



Deux idées : 
- adapter via des pipelines plus lourds mais contraintes d'infrastructures (Clariah Dane) le pipeline que nous avons proposé à l'aide de la computer vision nécessite un paramétrage pour chaque corpus pour le seuil de séquencage. et yolo ? comment l'adapter ? 


Proposer des outils plus performants que ceux proposés dans cette étude afin de pouvoir s’adapter à tous les types de corpus, tout en prenant en compte les contraintes d’infrastructures. Mais : contraintes d’infrastructures. Skinsoft héberge les données clients donc amener un algorithme lourd en local.  

Transformers : Revisiting Human-in-the-Loop Object Retrieval with Pre-Trained Vision Transformers  K. Zaher, O. Buisson et A. Joly ICPR 2026 - 28th International Conference on Pattern Recognition, 2026le 


En 2025, le ministère de la culture identifie "l'enrichissement de la connaissances et des savoirs" comme un usage potentiel de l'IA en secteur culturel avec par exemple "de l’aide au catalogage, à l’indexation automatique des collections, au tri et au classement des documents et données destinés à être archivés, à la transcription des textes manuscrits, à la reconnaissance de formes, à la diffusion et à la valorisation des contenus culturels et patrimoniaux quel que soit leur format (image, son…)". \footcite{2025a}